% !TEX root = root.tex
\documentclass[11pt, a4paper]{article} % Removed twocolumn
\usepackage[colorlinks=true, urlcolor=blue, linkcolor=red]{hyperref}
\usepackage{tabularx}
\usepackage{booktabs}
\usepackage[margin=1in]{geometry} % Standard margins for single column

\title{Predicting the likelihood of a TikTok video going viral using intrinsic and extrinsic signals}
\author{Oscar Rodriguez, Mahda Soltani, Ela Naz Sigin}
\date{March 10, 2026}

\begin{document}

\maketitle

\begin{abstract}
    Social media platforms, like TikTok have become one of the dominant
    channels through which information, entertainment, and political and
    cultural trends propagate at scale. A defining feature of these
    platforms is the extreme skewness in attention: while the vast majority
    of uploaded videos receive little to no exposure, a small fraction
    achieve \textit{viral} reach, accumulating orders of magnitude more
    views within a short time window. This heavy-tailed distribution of
    engagement naturally raises the question of how to operationalize and
    predict “virality” in a principled way.

    This investigation explores the use of Machine Learning techniques to
    pin down a formal framework for the task of predicting the likelihood of
    a TikTok video going viral using both intrinsic (e.g. video content) and
    extrinsic features (e.g. video author).
\end{abstract}

\section*{Introduction}
\dots

\section*{Methods}

This section goes over the different aspects of the Machine Learning pipeline
that we crafted for the given task, starting from data processing all the way
to the model architecture and reasoning behind each of the decisions.

\subsection*{Data collection and processing}

We started off with a base dataset called
\href{https://huggingface.co/datasets/The-data-company/TikTok-10M}{TikTok-10M}
from which we sourced a rich subset of examples that we complemented with a
scraping step to download the videos and publisher accounts to extract
additional signals (e.g. video bytes, aspect ratio, etc.) as well
as account information (e.g. account statistics, publishing history, etc.).

With all this raw data at hand, we produced a
\href{https://huggingface.co/datasets/rodmosc/viral}{curated dataset} for
our task. The following tables contain a detailed documentation of dataset
dimensions, encompassing intrinsic content signals and extrinsic author/metadata
signals.

\begin{table}[h]
\centering
\small
\begin{tabularx}{\textwidth}{|l l X l|}
    \hline
    \textbf{Dimension} & \textbf{Type} & \textbf{Description} & \textbf{Extraction / Treatment} \\
    \hline
    \multicolumn{4}{|l|}{\textbf{User / author Signals}} \\
    \hline
    user\_id/username & ID/String & Unique identifier and handle of the creator & Scraped (TikTok web profile) \\
    user\_language & String & Configuration language of the account & Base dataset \\
    city & String & City where the video was posted from & Base dataset \\
    user\_verified & Boolean & Verification status of the author & Scraped (TikTok web profile) \\
    is\_private & Boolean & Whether the account is set to private & Scraped (TikTok web profile) \\
    author\_stats & Float & Follower, heart, video, and friend counts & Scraped; $\log(1+x)$ scaled \\
    account\_create\_time & Timestamp & When the user account was first created & Scraped (unix epoch) \\
    \hline
\end{tabularx}
\end{table}

\begin{table}[h]
\centering
\small
\begin{tabularx}{\textwidth}{|l l X l|}
    \hline
    \multicolumn{4}{|l|}{\textbf{Video content \& metadata}} \\
    \hline
    description & String & User-generated caption and hashtags & Base dataset \\
    create\_time & Timestamp & Exact time of video publication & Base dataset (UTC) \\
    temporal\_features & Float & Hour of day and Day of week & Sine/cosine cyclic encoding \\
    is\_ad & Boolean & Flag for paid/sponsored content & Base dataset \\
    duration & Integer & Total length of the video in seconds & Base dataset \\
    challenges & String & Associated hashtag challenges or trends & Base dataset \\
    video\_bytes & Binary & Scraped video resized to 256x256 @ 1fps & \texttt{yt-dlp} + \texttt{moviepy} \\
    detected\_objects & String & Objects detected in video & pretrained facebook/detr-resnet-50 \\
    \hline
\end{tabularx}
\end{table}

\begin{table}[h]
\centering
\small
\begin{tabularx}{\textwidth}{|l l X l|}
    \hline
    \multicolumn{4}{|l|}{\textbf{Audio \& music specs}} \\
    \hline
    music\_info & String & Title, album, and author name of the audio & \texttt{yt-dlp} payload metadata \\
    music\_play\_url & String & URL source for the audio file & Base dataset \\
    \hline
\end{tabularx}
\end{table}

\begin{table}[h]
\centering
\small
\begin{tabularx}{\textwidth}{|l l X l|}
    \hline
    \multicolumn{4}{|l|}{\textbf{Technical / video specs}} \\
    \hline
    resolution & Integer & Video width and height in pixels & \texttt{yt-dlp} (Best format) \\
    aspect\_ratio & Float32 & Ratio of width to height & Calculated from payload \\
    filesize & Integer & Total storage size of the video file & \texttt{yt-dlp} payload metadata \\
    codecs & String & Video (vcodec) and audio (acodec) formats & \texttt{yt-dlp} payload metadata \\
    vq\_score & Float64 & Visual quality score & Base dataset \\
    \hline
\end{tabularx}
\end{table}

\begin{table}[h]
\centering
\small
\begin{tabularx}{\textwidth}{|l l X l|}
    \hline
    \multicolumn{4}{|l|}{\textbf{Engagement stats \& target labels}} \\
    \hline
    interaction\_counts & Integer & Plays, likes, comments, shares, and saves & Scraped + base dataset \\
    engagement\_score & Float64 & Weighted sum of interactions per view & Engineered; Weighted calculation \\
    view\_velocity & Float64 & Views accumulated per hour since creation & Engineered; $\log(1+x)$ scaled \\
    is\_viral & Binary & Classification ($P=0.95$) of engagement & Engineered (Quantile threshold) \\
    \hline
\end{tabularx}
\end{table}

\end{document}